% TOP_PROBLEM: {"id": 30006880, "title": "Existence of a p-Optimal Propositional Proof System", "category_id": 15, "category": {"id": 15, "name": "computer_science", "display_name": "Computer Science", "description": "Computational complexity, algorithms, and theoretical CS.", "slug": "computer-science", "order_index": 15, "created_at": "2026-07-31T15:26:25.671Z"}, "status": "open"}
% FIRST_POSTED: 2026-09-27
% !TeX program = lualatex
% ATTEMPT_STATUS: unresolved
\documentclass[11pt]{article}
\usepackage[margin=25mm]{geometry}
\usepackage{fontspec}
\setmainfont{Latin Modern Roman}
\usepackage{amsmath,amssymb,mathtools}
\usepackage{unicode-math}
\setmathfont{Latin Modern Math}
\usepackage{xurl,hyperref}
\hypersetup{colorlinks=true,urlcolor=blue}
\setcounter{secnumdepth}{0}
\setlength{\parindent}{0pt}
\setlength{\parskip}{5pt}
\setlength{\emergencystretch}{3em}
\begin{document}
\section*{\texorpdfstring{Existence of a p-Optimal Propositional Proof System}{Existence of a p-Optimal Propositional Proof System}}\label{30006880}
\subsection{Definitions and mathematical statement}
Let $\mathrm{TAUT}$ be the set of encoded propositional tautologies. A Cook--Reckhow proof system is a polynomial-time computable surjection $P:\{0,1\}^*\to\mathrm{TAUT}$; a string $\pi$ proves the formula $P(\pi)$. A system $P$ \emph{polynomially simulates} $Q$ if there is a polynomial-time function $t_Q$ such that
\[
 P(t_Q(\pi))=Q(\pi)\qquad\text{for all proof strings }\pi.
\]
A $p$-optimal system is one that polynomially simulates every Cook--Reckhow system. Determine whether such a $P$ exists. The translator and its polynomial bound may depend on $Q$, but $P$ must be fixed. Merely having polynomially short translations without an efficient procedure is the weaker notion of optimality and is not this target.
\par\smallskip\textit{Formulation and concept references:} \href{https://eccc.weizmann.ac.il/eccc-reports/1997/TR97-026/index.html}{[S1]}.

\subsection{Short English statement}
Is there one propositional proof system that can efficiently translate and simulate proofs from every other propositional proof system?
\subsection{Research attempt}
\paragraph{The target and current-source scope (27 September 2026).}
A translator must run in polynomial time and output a proof, not merely
certify that some short proof exists. Its polynomial can depend on the
simulated system. There is no demand for a common exponent over all
systems, and there is no promise that the verifier of the proposed
universal system can decide whether arbitrary programs are sound.
Those two distinctions drive the analysis below.

The supplied Messner--Tor\'an report [S1] defines the efficient simulation
question and states conditional existence results. A 2026 primary
preprint by Egidy [E1] studies relativization barriers for both positive
and negative answers. An oracle separation is not an unconditional
answer about the unrelativized set $\mathrm{TAUT}$. No theorem inspected
in this audit resolves the selected question.

This attempt constructs a universal system for finitely many competitors,
then for a suitably enumerable sound collection with explicit clocks.
It proves exactly why enumerating all candidate polynomial-time programs
fails, and separates the decision, short-proof, and proof-search routes.
All inputs and outputs are finite binary strings under fixed effective
encodings; malformed strings map to one fixed tautology $\tau$.

\paragraph{A concrete baseline system avoids a hidden existence assumption.}
Fix any polynomial-time syntax and evaluation algorithm for propositional
formulas. For a formula $\varphi$ with $v$ distinct variables, its truth
table can be checked in $2^v$ times a polynomial in $|\varphi|$.
Define a proof to contain $\varphi$ and unary padding of length $N$.
The baseline system $T$ parses this pair, checks that the truth-table
calculation fits a fixed polynomial budget in the total input length,
and, if so, executes it. It outputs $\varphi$ exactly when every
assignment satisfies it, and otherwise outputs $\tau$.

The budget can be implemented simply: run the evaluation loop for at
most $N$ elementary machine steps, with a fixed universal-simulation
polynomial overhead, and reject to $\tau$ if it has not finished.
Thus $T$ is polynomial-time on every string. It is sound by exhaustive
evaluation. For every tautology a sufficiently long padding allows its
finite truth-table calculation to finish, so $T$ is onto $\mathrm{TAUT}$.
No shortness guarantee is required for this baseline.

This construction illustrates a recurring principle: a potentially
expensive finite computation becomes polynomial in a \emph{padded proof}
length, but producing that padding need not be polynomial in the original
formula or competitor proof length. That second cost matters for
$p$-simulation.

\paragraph{Finite families have an explicit common upper bound.}
\textbf{Lemma 355.1.} Given finitely many Cook--Reckhow systems
$P_1,\ldots,P_s$, there is a system $R$ polynomially simulating all
of them.

\emph{Proof.} Encode a branch index $i$ followed by a proof $\pi$, and
put $R(i,\pi)=P_i(\pi)$ on valid branches, using $\tau$ otherwise.
A maximum of the finitely many runtime exponents is a constant, so
$R$ is polynomial-time. Its outputs are tautologies, and any branch
is already onto. The translator $\pi\mapsto(i,\pi)$ is linear-time
for each fixed $i$.\hfill$\square$

This construction also shows that simulation is a directed preorder:
identities give reflexivity, compositions of polynomial-time translators
give transitivity, and every finite set has an upper bound. A directed
preorder need not have a greatest element. For example the natural
numbers under their usual order have upper bounds for finite sets and
no largest element. Taking successively stronger finite combinations
therefore does not by itself produce a single polynomial-time limit.

The technical difficulty with infinitely many branches is visible even
before soundness: a branch with runtime $n^i$ cannot be executed by a
single fixed polynomial bound when $i$ is part of the input. Unary
clocks repair this timing problem; they do not repair unsound branches.

\paragraph{A uniform enumerable sound collection can be combined with clocks.}
Suppose an effective enumerator lists programs $Q_0,Q_1,\ldots$, each
of which is a Cook--Reckhow system. No computable common runtime exponent
is assumed. There is a single system $U$ that $p$-simulates every listed
system.

A proof for $U$ specifies an index $i$, a string $\pi$, and a unary
clock of length $N$. Run the enumerator for at most $N$ steps to obtain
its $i$th program. If it has not appeared, output $\tau$. If it has,
simulate that program on $\pi$ for at most $N$ steps and output the
result if the simulation terminates, otherwise output $\tau$.
One can include the baseline $T$ as an additional branch to guarantee
surjectivity even if the enumerated collection is empty.

This verifier has polynomial runtime in its total input length, with
fixed exponent: the simulations are clocked, the program description
produced in $N$ steps has length at most $N$, and ordinary universal
simulation has polynomial overhead. Soundness follows from the assumed
soundness of every enumerated program, not from a test performed by $U$.

For a fixed $i$, choose constants bounding the time until its description
is enumerated and its polynomial runtime on length $n$. The translator
outputs $(i,\pi,1^{N_i(|\pi|)})$, where $N_i$ is a polynomial exceeding
both bounds. Printing the unary clock is polynomial-time for this fixed
$i$. This proves the simulation with precisely the allowed dependence
of the polynomial on the competitor. The construction is fully uniform
as a verifier, even though its existence proof uses different clock
polynomials for different listed systems.

\paragraph{An exact cofinal-enumeration formulation.}
Call a family of systems $p$-cofinal if every Cook--Reckhow system is
$p$-simulated by some member of that family. The preceding argument and
transitivity prove the following equivalent formulation:
\[
 \text{a $p$-optimal system exists}
 \quad\Longleftrightarrow\quad
 \text{there is an effectively enumerable $p$-cofinal family of sound systems}.
 \tag{355.1}
\]
For the reverse implication combine the enumerable family with clocks.
For the forward implication the singleton consisting of a $p$-optimal
system is already an effectively enumerable cofinal family. The latter
observation is logically enough and does not purport to recognize that
singleton effectively from the property of optimality.

One can also exhibit a richer enumerable family once a fixed sound
system $P$ is given. Enumerate all clocked polynomial-time functions
$f_e$, and define $Q_e$ to run $P(f_e(\pi))$ on one branch and $P(\pi)$
on another. Each $Q_e$ is a sound surjection. If $P$ is $p$-optimal,
then translators into $P$ are represented among these clocks. This
construction never has to decide the soundness of an arbitrary external
program, because all outputs are routed through an already sound $P$.
It does not construct that $P$ in the first place.

\paragraph{Why enumeration of all polynomial-time programs is unsound.}
The syntax of polynomial-time programs can be effectively enumerated
by supplying clocks. Their semantic property of always outputting
tautologies cannot be enumerated in the same way. Here is a direct
reduction demonstrating the obstruction, even when surjectivity is
built into the candidate.

Given a Turing machine index $e$, construct a clocked polynomial-time
function $F_e$ with two input branches. On branch zero, run the baseline
system $T$. On branch one with payload length $n$, simulate machine
$e$ on its blank input for $n$ steps. If it halts within that time,
output a fixed contradiction $\bot$; otherwise output $\tau$.
For each fixed $e$ this is a polynomial-time function, with a clock
that can be supplied effectively. It is onto the tautologies through
its baseline branch regardless of $e$.

If $e$ never halts, every output is a tautology, so $F_e$ is a valid
proof system. If $e$ halts, a long enough branch-one input outputs
$\bot$, so $F_e$ is not a proof system. Hence
\[
 F_e\text{ is a Cook--Reckhow system}
   \quad\Longleftrightarrow\quad e\text{ does not halt}.
 \tag{355.2}
\]
The set of nonhalting machine indices is not recursively enumerable:
if both halting and nonhalting were enumerable, dovetailing the two
lists would decide the halting problem, contradicting its diagonal
undecidability proof. Thus the set of \emph{all program descriptions}
that compute Cook--Reckhow systems is not recursively enumerable.

This is not a proof that (355.1) fails. A cofinal enumerable collection
need not contain every correct program description, or even compute
every proof-system function. It only needs enough sound systems to
simulate each competitor. Confusing exhaustive semantic recognition
with cofinality would turn a valid undecidability observation into an
invalid disproof of $p$-optimality.

\paragraph{Truth-table guarding repairs soundness but loses the translator.}
A universal verifier could execute a submitted clocked program on its
submitted proof, obtain a formula, and check that formula by a padded
truth table before outputting it. This is a sound construction and can
be made polynomial in the total padded input size. But a formula with
$v$ variables may require $2^v$ evaluations. A competitor may output
such a formula from a proof of length polynomial in $v$.

The proposed translator would then have to print exponentially many
padding symbols. The fact that the target verifier is polynomial in
its own proof length says nothing about this translation cost. One
cannot remove the guard merely because the input claims that its
program is sound: false assertions of soundness are still legal
binary strings on which a Cook--Reckhow function must output a tautology.

A finite soundness check is insufficient for the same reason as (355.2).
The program there can agree with a valid system on every proof length
up to an arbitrary bound, then output a contradiction when the simulated
machine eventually halts. Testing finitely many proof strings cannot
establish the universal soundness condition.

\paragraph{A complete conditional construction under $P=NP$.}
If $P=NP$, then $P=\mathrm{coNP}$ and tautology testing is polynomial-time.
Define $P_*(x)$ to output $x$ if it is a well-formed tautology and $\tau$
otherwise. This is a polynomial-time surjection onto $\mathrm{TAUT}$.
For any proof system $Q$, the translator
\[
 t_Q(\pi)=Q(\pi)
\]
is polynomial-time and satisfies $P_*(t_Q(\pi))=Q(\pi)$.
Thus $P=NP$ implies existence of a $p$-optimal system by a particularly
transparent construction. It is a conditional implication, not an
assumption available in the selected problem.

More generally, it suffices to have a fixed sound system $P$ and a
polynomial-time algorithm which, on every tautology $\varphi$, finds
a $P$-proof of $\varphi$ with polynomial output length. Then composing
that algorithm with $Q$ produces the required translator. An algorithm
promised to work only on tautologies is enough for this composition,
because every $Q(\pi)$ is a tautology. Such a proof-search hypothesis
is additional substantive information, not a consequence of knowing
that short proofs exist.

\paragraph{Short proofs yield weaker optimality, not an efficient translation.}
Call $P$ polynomially bounded when every tautology $\varphi$ has some
$P$-proof of length at most a fixed polynomial in $|\varphi|$.
Then for every $Q$, polynomially short translations exist as strings:
$|Q(\pi)|$ is polynomial in $|\pi|$, since $Q$ must print its output,
and the promised $P$-proof length is consequently polynomial in $|\pi|$.
This establishes simulation in the proof-length sense only. Exhaustively
searching all possible short $P$-proofs takes exponential time in their
length bound. No efficient choice function follows from the existential
length bound.

For completeness, the standard equivalence with $NP=\mathrm{coNP}$
can be verified directly. A polynomially bounded system gives an NP
witness relation for tautology: guess the short proof and evaluate $P$.
Since $\mathrm{TAUT}$ is coNP-complete, this yields $NP=\mathrm{coNP}$.
Conversely, if tautology has an NP verifier $V(\varphi,z)$ with
polynomially bounded certificates, define a system to output $\varphi$
on accepted pairs and $\tau$ otherwise. It is sound, surjective, and
polynomially bounded. The standard NP/coNP completeness facts are the
only complexity-theoretic inputs in this equivalence.

Accordingly, this argument from $NP=\mathrm{coNP}$ supplies a
length-optimal system, but it does not on its own supply the
polynomial-time translators required here. It also does not show that
$p$-optimal systems must be polynomially bounded. A universal simulator
may still give very long shortest proofs for some tautologies.

\paragraph{Padding makes the last distinction especially concrete.}
Starting from any system $P$, define a padded system $P^{\rm pad}$
whose nondefault proofs encode $(\pi,1^{2^{|\pi|}})$ and whose output
is $P(\pi)$. The verifier can check the padding length by comparing
its binary length counter with $2^{|\pi|}$, and can run $P$ within a
fixed polynomial in the padded length. It remains a sound surjection.
A direct map from $\pi$ to that particular padded proof is exponential,
although deleting the padding gives an efficient translator in the
opposite direction.

This example does not prove that no other short padded proof of the
same formula exists; $P$ might have alternative proofs. Its limited
purpose is to show concretely why a verifier's polynomial runtime and
an explicit string-level translation are different requirements.
A purported lower bound on all translations would need a proof-length
lower bound for every alternative proof, which is not provided here.

\paragraph{Adversarial checks and first gap.}
The finite diagnostic implements the truth-table baseline on small
formulas, finite branch combinations and a delayed-unsoundness program.
It checks that every accepted baseline output is a tautology, that branch
translations preserve outputs, and that a guard tested only below the
delay misses the later contradiction. These are exact Boolean checks,
not an experiment purporting to decide a universal semantic property.

\textbf{Outcome: UNRESOLVED.} The construction for finite families and
sound enumerable cofinal collections is complete, with simulation clocks
and output lengths accounted for. The exact remaining gap is to obtain
such a cofinal collection, or another universal sound verifier with
efficient translators, without deciding the unenumerable semantic set
in (355.2) and without paying truth-table padding costs. The conditional
$P=NP$ proof, weak short-proof simulation, and 2026 oracle results do
not close that gap unconditionally.

\subsection{Sources}
\begin{itemize}
\item[S1]Optimal proof systems for propositional logic and complete sets.
\url{https://eccc.weizmann.ac.il/eccc-reports/1997/TR97-026/index.html}.
\textit{Formulation source.}
\item[E1]Fabian Egidy, \emph{The SPARSE-Relativization Framework and
Applications to Optimal Proof Systems}, arXiv:2602.02294 (2026).
\url{https://arxiv.org/abs/2602.02294}.
\textit{Primary abstract inspected 27 September 2026; oracle and
relativization results are distinguished from an unrelativized resolution.}
\end{itemize}
\end{document}
